% Source: https://artofproblemsolving.com/community/c3h1743708_problem_5__perpendicular_intersections
% Stable identity: aops-topic-1743708-q01
The line $2x+y=3$ has slope $-2$, so a perpendicular line has slope $\dfrac{1}{2}$. With $y$-intercept $3$ it is $y=\dfrac{1}{2}x+3$, which meets the $x$-axis where $\dfrac{1}{2}x=-3$, at $(-6,0)$. The answer is $\boxed{\textbf{(D)}\ (-6,0)}$.

\begin{asy}
usepackage("amsmath");
size(6cm);
defaultpen(fontsize(10pt));
import graph;
real f(real x){return 3-2x;} real g(real x){return 3+x/2;}
draw((-8,0)--(4,0),gray,Arrow(5));draw((0,-2)--(0,7),gray,Arrow(5));
draw((-0.5,4)--(2.5,-2),linewidth(0.9));draw((-8,-1)--(3,4.5),blue+linewidth(0.9));
pair X=(-6,0),Y=(0,3),Z=(0.6,3.3);
dot(X);dot(Y);
label("$x$",(4,0),S);label("$y$",(0,7),W);
label("$(0,3)$",Y,2*dir(-14));label("$(-6,0)$",X,dir(-60));
label("$2x+y=3$",(2.2,-1.4),E);label("$y=\dfrac{1}{2}x+3$",(-5.6,3.2),blue);
pair H=(0.6,1.8);
label("\mbox{}",(min(currentpicture,user=true)+max(currentpicture,user=true))/2,black); // keeps the text after the figure black when LaTeX typesets the labels inline
\end{asy}
